% Last updated: 2025/04/23

The primary data source used is the Medicare Physician and Other Practitioners by Provider and Service Dataset, also known as the Public Use Files (PUF). The PUF reports utilization and payment data each year for 2013-2021, aggregated by physician and service performed. Services are identified by Healthcare Common Procedure Coding System (HCPCS) code. HCPCS codes are standardized procedure codes used in the billing process to describe services performed by a physician during a visit. For example, a single observation of this dataset could describe the number of lower back laminectomies a physician performs in a year. This physician-service-year level dataset is based on Medicare claims data for all beneficiaries enrolled in the part B fee-for-service program. It includes all physicians who submitted part B non-institutional claims, which is almost the universe of all practicing physicians. Services performed by physicians who bill as part of an organization cannot be identified separately from their organization. Any service code that a physician bills for less than 11 times during the year is not reported in the PUF, which may lead to underestimates of number of surgeries performed.\footnote{To avoid including surgeons who may be observed performing no surgeries each year due to HCPCS code level censoring, I drop any surgeon for whom there is a large share of administratively reported total services that are unassociated with services codes categorized as surgery, non-surgical treatment, diagnostics, and dependent. This allows me to keep surgeons who likely genuinely have low surgical volume in my sample while dropping those who may have larger surgical volume that is censored.}

I restrict my sample of physicians to MD's with a surgical subspecialty recognized by the American College of Surgeons, excluding pediatric surgery and oral and maxillofacial surgery.\footnote{The ACS recognizes cardiothoracic surgery, colorectal surgery, general, obstetrics and gynecology, gynecologic oncology, neurosurgery, ophthalmic surgery, oral and maxillofacial surgery, orthopedic surgery, otolaryngology (ENT), pediatric surgery, plastic surgery, urology, and vascular surgery as surgical subspecialties. Cardiothoracic surgery is broken down into cardiac surgery and thoracic surgery in my data.} Pediatric surgery is not relevant for the Medicare population, and oral surgeons, unlike other surgeons, attend dental school instead of medical school. Since I link information in medical school hospital referral region (HRR) at the time of graduation to currently practicing surgeons, I only look at surgeons who graduate medical school in or after 2013.\footnote{I additionally drop surgeons for whom graduation year is likely listed incorrectly in CMS Physician Compare based on the length of GME required by their listed surgical subspecialty, with 1 year of leeway.}

\subsection{Treatment Intensity}
My main outcome of interest is physician treatment intensity. This measure is calculated as the number of surgeries a surgeon performs divided by the number of Medicare patients seen in a particular year. In order to determine service codes pertaining to surgeries, I evaluate the description of each HCPCS code, then manually assign each code to one of four categories: surgery, diagnostic procedure, office visit, or ``dependent" service. I consider a dependent service to be those often billed together with a larger procedure; for example, anesthesia is often billed together with surgery and is not a stand-alone procedure.  Using this classification system, I compute treatment intensity at the physician-year level.

My main independent variable is the average treatment intensity among all physicians practicing within a graduating physician's medical school HRR at the time they graduate from medical school. Treatment intensity is calculated as described above. The distribution of treatment intensity across HRRs during the time my sample graduates from medical school is show in Figure \ref{fig:surg-hrr}. I link medical school HRR intensity using each surgeon's medical school and graduating year, which are from the CMS Physician Compare National Downloadable File. I remove any physician observations listing a medical school that closed prior to 2013, as well as any observations whose medical school is missing or labeled as ``OTHER." The majority of the latter group is comprised of physicians who completed medical school at an institution outside of the US. I collect zip codes for medical schools from their official websites. These zip codes are used to assign each medical school to a HRR using crosswalk files from the Dartmouth Atlas.

\begin{figure}
     \centering
     \begin{subfigure}[b]{0.45\textwidth}
         \centering
         \includegraphics[width=\textwidth]{surg2013.png}
         \label{fig:surg2013}
     \end{subfigure}
     \hfill
     \begin{subfigure}[b]{0.45\textwidth}
         \centering
         \includegraphics[width=\textwidth]{surg2018.png}
         \label{fig:surg2018}
     \end{subfigure}
    \caption{Treatment intensity by HRR}
    \label{fig:surg-hrr}
\end{figure}

\subsection{Physician Location}
Physician zip codes are identified from Medicare Data on Provider Practice and Specialty (MD-PPAS) for 2013-2018 and from Physician Compare for 2019-2022. Physician practice addresses can also be found in the National Plan and Provider Enumeration System (NPPES) for this time period. However, location in NPPES is only updated when a physician chooses to update it, and this information may be inaccurate or outdated as a result. Addresses from MD-PPAS and Physician Compare are aggregated from Medicare claims, which allows me to determine physician location with greater certainty. From these addresses, I use zip code to assign each physician-year observation to their current practice HRR. 

There are a number of physicians who bill from multiple zip codes. This is not a problem for location reported in MD-PPAS, since it reports the primary zip code of a practitioner based on the number of services performed. The information needed to create a similar primary zip code variable for Physician Compare cannot be linked to each practice location. Within the Physician Compare data, I only consider physicians practicing in one core-based statistical area (CBSA) since these physicians likely have a single, primary practice area. I then assign HRR based on zip code and only keep observations with both a medical school HRR and a single current practice HRR. I define physician movers as those who practice in a HRR different from the HRR where they completed medical school.

\subsection{Other Included Variables}
I also control for physician characteristics and aggregate Medicare patient characteristics for all beneficiaries seen by the physician that year. Physician characteristics include gender, graduation year, and indicators for surgical subspecialty. Aggregate Medicare patient characteristics include average age, proportion female, proportion Black, proportion dual-enrolled in Medicaid, and average hierarchical condition category (HCC) risk score. Dual-enrollment in Medicaid serves as a proxy for the proportion of the physician's patient population that is low-income. HCC risk score is a risk-adjusted estimate of a patient's healthcare costs, and a higher score is indicative of higher expected spending. I use this as a proxy for case-mix, since patients whose cases are more severe may require more surgical treatment. All physician and patient characteristics are from in the PUF, except for graduation year, which is from Physician Compare. I include indicator variables for missing patient characteristics to maintain sample size.

\newcommand*{\tabindent}{ \hspace{3mm}}

\setlength{\LTpost}{0mm}
\LTcapwidth=\textwidth
\begin{ThreePartTable}
\begin{TableNotes}[para, flushleft]
    \textit{Notes.} Missing patient characteristics imputed with unconditional mean among sample to retain sample size.
\end{TableNotes}
\begin{longtable}{lccccc}
    \label{tab:summ} \\
    \caption{Summary statistics} \\
    \toprule
    \multicolumn{1}{l}{Variable} & Mean & Std. Dev & Min & Max & N \\ 
    \midrule\addlinespace[2.5pt]
    \multicolumn{6}{l}{\textbf{Treatment Intensity}} \\ 
    \midrule\addlinespace[2.5pt]
    \tabindent Physician intensity & $0.21$ & $0.45$ & $0$ & $17.977$ & $6,760$ \\ 
    \tabindent Med school HRR intensity & $0.22$ & $0.031$ & $0.122$ & $0.423$ & $6,760$ \\ 
    %\tabindent Med school teaching hospital intensity & - & - & - & - & - \\ 
    \multicolumn{6}{l}{Physician Characteristics} \\ 
    \midrule\addlinespace[2.5pt]
    \tabindent Mover & $0.854$ & $0.354$ & $0$ & $1$ & $6,760$ \\ 
    \tabindent Female & $0.305$ & $0.46$ & $0$ & $1$ & $6,760$ \\ 
    \tabindent Graduation year & $2013.928$ & $1.05$ & $2013$ & $2018$ & $6,760$ \\ 
    \tabindent Cardiac surgeon & $0.003$ & $0.05$ & $0$ & $1$ & $6,760$ \\ 
    \tabindent Colorectal surgeon & $0.002$ & $0.044$ & $0$ & $1$ & $6,760$ \\ 
    \tabindent General surgeon & $0.107$ & $0.309$ & $0$ & $1$ & $6,760$ \\ 
    \tabindent Gynecologic oncologist & $0.001$ & $0.027$ & $0$ & $1$ & $6,760$ \\ 
    \tabindent Hand surgeon & $0.008$ & $0.089$ & $0$ & $1$ & $6,760$ \\ 
    \tabindent Neurosurgeon & $0.011$ & $0.105$ & $0$ & $1$ & $6,760$ \\ 
    \tabindent OBGYN & $0.051$ & $0.22$ & $0$ & $1$ & $6,760$ \\ 
    \tabindent Ophthalmologist & $0.37$ & $0.483$ & $0$ & $1$ & $6,760$ \\ 
    \tabindent Orthopedic surgeon & $0.201$ & $0.401$ & $0$ & $1$ & $6,760$ \\ 
    \tabindent Otolaryngologist (ENT) & $0.114$ & $0.318$ & $0$ & $1$ & $6,760$ \\ 
    \tabindent Plastic surgeon & $0.003$ & $0.058$ & $0$ & $1$ & $6,760$ \\ 
    \tabindent Surgical oncologist & $0.002$ & $0.042$ & $0$ & $1$ & $6,760$ \\ 
    \tabindent Thoracic surgeon & $0.004$ & $0.064$ & $0$ & $1$ & $6,760$ \\ 
    \tabindent Urologist & $0.108$ & $0.31$ & $0$ & $1$ & $6,760$ \\ 
    \tabindent Vascular surgeon & $0.014$ & $0.117$ & $0$ & $1$ & $6,760$ \\ 
    \midrule\addlinespace[2.5pt]
    \multicolumn{6}{l}{Patient Characteristics} \\ 
    \midrule\addlinespace[2.5pt]
    \tabindent Average age & $72.921$ & $3.009$ & $43$ & $84$ & $6,760$ \\ 
    \tabindent Proportion female & $0.568$ & $0.148$ & $0$ & $1$ & $6,760$ \\ 
    \tabindent Proportion Black & $0.142$ & $0.087$ & $0$ & $0.887$ & $6,760$ \\ 
    \tabindent Proportion dual-enrolled & $0.196$ & $0.128$ & $0$ & $1$ & $6,760$ \\ 
    \tabindent Average HCC risk score & $1.445$ & $0.728$ & $0.53$ & $11.413$ & $6,760$ \\ 
    \tabindent Proportion female missing & $0.023$ & $0.15$ & $0$ & $1$ & $6,760$ \\ 
    \tabindent Proportion Black missing & $0.625$ & $0.484$ & $0$ & $1$ & $6,760$ \\ 
    \tabindent Proportion dual-enrolled missing & $0.155$ & $0.362$ & $0$ & $1$ & $6,760$ \\ 
    \bottomrule
    \insertTableNotes
\end{longtable}
\end{ThreePartTable}


\subsection{Sample Restrictions}
All treatment intensity measures are based on surgeries performed on Medicare beneficiaries with part B fee-for-service coverage. These patients are older than the general population and could experience more harm from surgery; thus, treatment intensity may be lower for this patient group compared to a more general one. Further, fee-for-service beneficiaries are likely less healthy than those in Medicare Advantage due to selection. Due to age and other risk factors, it is probable that physicians are less likely to deviate from treatment norms either from their training or from their current workplace. That is, the Medicare population is less healthy compared to the general population, and physicians may tend to choose treatment decisions that are more ``by the book" and more impacted by medical school training. As a result, the magnitude of any result found here can be thought of as an upper bound for the general patient population.
